\documentclass[main.tex]{subfiles}


\begin{document}

%from pagenumber to up
% \phantomsection 
% \hypertarget{toc}{}

 

 

 
   

\section{Катушка. Самоиндукция}


% Опыт показывает, что магнитный поток через контур, создаваемый током в этом контуре, пропорционален силе этого тока: $\Phi\sim I$.  
% \textbf{Индуктивность} ($L$ [Гн])~--- это характеристика катушки\footnote{Более точно, контура.}, показывающая ее способность создавать магнитное поле:
% \begin{equation}
%     L=\frac{\Phi_L}{I},
% \end{equation}
% где $\Phi_L$~--- магнитный поток через катушку, создаваемый током $I$ в катушке.

\textbf{Магнитный поток катушки}, создаваемый током в ней и пронизывающий ее, равен:
\begin{equation}
    \Phi_L=LI,
\end{equation}
где $L$~--- \emph{индуктивность} катушки\footnote{Индуктивность~--- это характеристика катушки (контура), показывающая ее способность создавать магнитное поле. В общем случае индуктивность катушки определяется так: $L=\dfrac{\mu\mu_0N^2S}{l}$, где $\mu$~--- магнитная проницаемость вещества внутри катушки, $\mu_0$~--- магнитная постоянная, $N$~--- число витков катушки, $S$~--- площадь поперечного сечения катушки, $l$~--- длина катушки.}, $I$~--- ток в катушке.


% Например, из двух катушек одинаковых размеров с равными токами внутри себя создает более сильное поле та катушка, у которой индуктивность больше. 

Говорят, что индуктивность~--- это способность катушки накапливать энергию магнитного поля.

\textbf{Энергия катушки} рассчитывается по формуле:
\begin{equation}
    W_L=\frac{LI^2}{2}.
\end{equation}

\textbf{Самоиндукция}~--- это явление возникновения ЭДС в контуре (катушке) при изменении силы тока в этом же самом контуре (катушке). Схема опыта по обнаружению явления самоиндукции показана на рис.\,\ref{fig:p164} (слева).



    \begin{figure}[H]
    \centering
\begin{circuitikz}[circuit ee IEC,font=\small,thin,
    line cap=round,line join=round,
>={Stealth[inset=0pt,round,length=3mm, angle'=20]},
vec/.style={>={Stealth[inset=0pt,round,length=3.5mm, angle'=20]}}]

    \ctikzset{bipoles/americaninductor/coils=3}


\draw (0,0) to[*battery1,l={$\mathscr{E},r$},invert,name=E] ++ (2.5,0) to[make contact={info={[black]$K$}}] ++ (2.5,0) --++ (0,-2.5) --++ (-.4,0) coordinate (l);
\draw[line cap=butt] (l) to[*L=,l_=$L$,mirror] ++ (-2.1,0) coordinate (ll);
\draw (ll) to[*lamp,l_=Л] ++ (-2.5,0) --++ (0,2.5);


%%%%
%%%%
% right pic
\begin{scope}[xshift=7cm]
    
\draw (0,0) to[*battery1,l={$\mathscr{E},r$},invert,name=E] ++ (2.5,0) --++ (1.25,0) node[above] {$K$} --++ (1.25,0) --++ (0,-2.5) --++ (-.4,0) coordinate (l);
\draw[line cap=butt] (l) to[*L=,l_=$L$,mirror] ++ (-2.1,0) coordinate (ll);
\draw (ll) to[*lamp,l_=Л] ++ (-2.5,0) --++ (0,2.5);

\def\sl{7.5*\pgflinewidth}
\def\bl{15*\pgflinewidth}
\def\di{2.5*\pgflinewidth}
\draw[shift={(4.5,-2.5)},gray,thick] (0,0) ++ (-\di,\sl) --++ (0,{-2*\sl}) ++ ({2*\di},{\sl}) ++ (0,\bl) --++ (0,{-2*\bl}) ++ (-\di,{2*\bl}) 
% node[fill,circle,minimum size=1mm,inner sep=0] {}
node[above,black] {$\mathscr E_{si}$}
;
\end{scope}


%\Longrightarrow
\path (6,-1.25) node {$\Longrightarrow$};


\end{circuitikz}
\caption{К опыту по обнаружению явления самоиндукции}
\label{fig:p164}
\end{figure}

Имеется цепь, состоящая из последовательно соединенных катушки индуктивности~$L$ и лампы~Л, подключенных к батарейке~$\mathscr E,r$ через ключ~$K$. При разомкнутом ключе тока в цепи нет, лампа не горит. 

Пусть теперь ключ в рассмотренной схеме замыкают (рис.\,\ref{fig:p164}, справа).

%     \begin{figure}[H]
%     \centering
% \begin{circuitikz}[circuit ee IEC,font=\small,thin,
%     line cap=round,line join=round,
% >={Stealth[inset=0pt,round,length=3mm, angle'=20]},
% vec/.style={>={Stealth[inset=0pt,round,length=3.5mm, angle'=20]}}]

%     \ctikzset{bipoles/americaninductor/coils=3}

% %%%%
% %%%%
% % right pic

% \draw (0,0) to[*battery1,l={$\mathscr{E},r$},invert,name=E] ++ (2.5,0) --++ (1.25,0) node[above] {$K$} --++ (1.25,0) --++ (0,-2.5) --++ (-.4,0) to[*L=,l_=$L$,mirror] ++ (-2.1,0) to[*lamp,l_=Л] ++ (-2.5,0) --++ (0,2.5);

% \def\sl{7.5*\pgflinewidth}
% \def\bl{15*\pgflinewidth}
% \def\di{2.5*\pgflinewidth}
% \draw[shift={(4.5,-2.5)},gray,thick] (0,0) ++ (-\di,\sl) --++ (0,{-2*\sl}) ++ ({2*\di},{\sl}) ++ (0,\bl) --++ (0,{-2*\bl}) ++ (-\di,{2*\bl}) 
% % node[fill,circle,minimum size=1mm,inner sep=0] {}
% node[above,black] {$\mathscr E_{si}$}
% ;


% \end{circuitikz}
% \caption{Схема к опыту (самоиндукция): сразу после замыкания ключа}
% \label{fig:p165}
% \end{figure}

\emph{Сразу после} замыкания ключа~$K$ тока в цепи \emph{нет.} Это объясняется тем, что в катушке возникает \emph{ЭДС самоиндукции~$\mathscr E_{si}$, препятствующая изменению тока через катушку}\footnote{Говорят, что ток в катушке не может измениться \emph{скачком.}} (постоянный ток в цепи устанавливается спустя некоторое время~--- соответственно, лампа загорается \emph{постепенно}).

\textbf{ЭДС самоиндукции} в катушке при изменении тока в ней находится по формуле:
\begin{equation}
    \mathscr E_{si}=-L\frac{\Delta I}{\Delta t}=-LI'.
\end{equation}

ЭДС индукции будет возникать в катушке при любом способе изменения магнитного потока катушки (по закону электромагнитной индукции Фарадея): $\mathscr E_{i}=-\dfrac{\Delta \Phi_L\strut}{\Delta t\strut}=-\Phi_L'$.






 
\end{document}